\documentclass{article}
\usepackage[utf8]{inputenc}
\usepackage[T1]{fontenc}
\usepackage{geometry}
\geometry{
    a4paper,
    left=2.5cm,
    right=2.5cm,
    top=2.5cm,
    bottom=2.5cm
}
\usepackage{graphicx}
\usepackage{float}
\usepackage{hyperref}
\usepackage[english]{babel}
\usepackage{amsmath}
\usepackage{amssymb}
\usepackage{array}
\usepackage{booktabs}
\usepackage{caption}
\usepackage{subcaption}
\usepackage{enumitem}
\usepackage{fancyhdr}
\usepackage{titlesec}
\usepackage{url}
\usepackage{textcomp}
\usepackage{xcolor}
\usepackage{listings}
\usepackage{multirow}
\usepackage{longtable}
\usepackage{tabularx}
\usepackage{times}

% Set up page style
\pagestyle{fancy}
\fancyhf{}
\rfoot{\thepage}
\renewcommand{\headrulewidth}{0pt}

% Title formatting
\titleformat{\chapter}[display]{\normalfont\bfseries}{\chaptertitle\thechapter}{0pt}{\Huge}
\titlespacing*{\chapter}{0pt}{50pt}{40pt}

% Document info
\title{Template Document}
\author{Academic Support Agent}
\date{\today}

\begin{document}
\maketitle

\includepdf{cover.pdf} % .
	
    \renewcommand\thetable{\thesection.\arabic{table}} 
    \renewcommand\thefigure{\thesection.\arabic{figure}} % 
    
    \newCJKfontfamily{\heiti}[AutoFakeBold={1.7}]{SimHei.TTF}
    %\newCJKfontfamily{\songti}[AutoFakeBold={3}]{SimSun.TTC}

	
    \renewcommand{\figurename}{ }
    \renewcommand\tablename{ }% table 
    \captionsetup{labelformat=default,labelsep=default}% 
    \renewcommand{\thesubfigure}{\empty}% subfigure a,b,c 


    %%  
    %
    \fancyhead[C]{\zihao{5}\songti{ }}
    \phantomsection
    \addcontentsline{toc}{section}{ }

    
    \vspace*{0pt}
    \centerline{\zihao{3}\heiti\bfseries  \quad  }
    \vspace*{16pt}

    \linespread{1.4}
    \zihao{-4} 

     
    
    
    
    
    
    
    
    
    
    
    
    
    
    
    
    
    
    
    
    
    
    
    
    
    
    
    



    \bigskip

    \noindent{\zihao{-4}\songti \textbf{ :}}
     
    \pagenumbering{Roman}% 



    \newpage
    \phantomsection\addcontentsline{toc}{section}{Abstract}% 
    \fancyhead[C]{\zihao{5}\songti{Abstract}}
    \vspace*{0pt}
    \centerline{\zihao{3}\bfseries Abstract}
    \vspace*{20pt}

    \linespread{1.4}\zihao{-4}
    *************** ****************
   

    \bigskip
    \noindent\textbf{\zihao{-4} Keywords:} 
    believe in yourself always



    \newpage
    \setcounter{secnumdepth}{4}% 
    \setcounter{tocdepth}{2}% 2 

    \renewcommand\contentsname{\zihao{3}\heiti\textbf{ \quad  }}
  
    \begin{spacing}{1.7}
        {\centering
    \end{spacing}
    \thispagestyle{mulu}
    \vspace*{16pt}
	




    \setcounter{page}{0} % 0 
    %\thispagestyle{empty}% 

    \pagestyle{yemei}
    \newpage
    \pagenumbering{arabic}% 
    \setcounter{page}{1} % 1 

    \section{ }
    % \section{} 
    \markboth{\text{ }\thesection\text{ } \quad  }{\text{ }\thesection\text{ } \quad  }   

    \subsection{ }
    
         
       ; 
      
      
       
       ; 
      
      
    


    

    \subsubsection{ }
    $\times \times \times\times \times \times\times \times \times\times \times \times\times \times \times\times \times \times\times \times \times$

    \subsubsection{$\times \times \times\times \times \times$}
    $\times \times \times\times \times \times\times \times \times\times \times \times\times \times \times\times \times \times\times \times \times$

    \subsection{ }

    





    \newpage
    \section{ }
    % \section{} 
    \markboth{\text{ }\thesection\text{ } \quad  }{\text{ }\thesection\text{ } \quad  }  


    \subsection{ }
    %%%% $$ $$$$
     $a^2+b^2=c^2$.% 

     $$a^2+b^2=c^2.$$% 

     math % math 
    \begin{math}
        a^2+b^2=c^2.
    \end{math}

    %%%% equation 
    equation 
    \begin{equation}% 
        x^3+x^2+x+1=0
    \end{equation}

    equation* 
    \begin{equation*}% \tag{} 
        x^3+x^2+x+1=0
        \tag{2.1}
    \end{equation*}

    \subsection{ }
     equation 
    \begin{equation}
    \label{f1}% 
        \left\{ % 
        \begin{aligned}
            &\phi _t =-(\phi^3 -g\phi ^2-\epsilon \phi +(1+\Delta  )^2 \phi),\\
            &\phi(x,y,0)=\phi_0(x,y),\\
            &\phi (x+L,y,t)=\phi(x,y,t),\\
            &\phi (x,y+L,t)=\phi (x,y,t),
        \end{aligned}
		\right.\tag{2.2}
    \end{equation}

     ref $(\ref{f1})$ 

    \subsection{ }
     cite $\cite{ref1}$ pdf 









    \newpage
    \section{ }
     

    % \section{} 
    \markboth{\text{ }\thesection\text{ } \quad  }{\text{ }\thesection\text{ } \quad  }  


    \subsection{ }
    LaTex tikz. 

     main.tex 

     

    \begin{figure}[htbp]
        \centering
        \includegraphics[height=6cm,width=7.5cm]{0.jpg} 
        \caption{\songti{ .}}
        \captionsetup{font={small,bf,stretch=1}}
        \label{figure1}
    \end{figure}

     
    \begin{figure}[htbp]
        \centering
        \subfigure[$T=0$]{
        \includegraphics[height=4cm,width=4.5cm]{0.jpg} 
        }
        \quad
        \subfigure[$T=10$]{
        \includegraphics[height=4cm,width=4.5cm]{10.jpg} 
        }
        \quad
        \subfigure[$T=20$]{
        \includegraphics[height=4cm,width=4.5cm]{20.jpg} 
        }
        \caption{\songti{ }}
        \captionsetup{font={small,bf,stretch=1}}
        \label{figure2}
    \end{figure}

     ref $\ref{figure1}$ $\ref{figure2}$ 
    

    \subsection{ }
     latex 3 

     1 
    \begin{table}[htbp]
        \centering
        \setlength{\tabcolsep}{10mm}
        \begin{tabular}{ccc}  
            \hline
             & ( )& \\ \hline
            $1900$&$10$&--\\ 
            $2000$&$100$&$1000\%$\\ 
            $2100$&$1000$&$1000\%$\\  \hline
        \end{tabular} 
        \caption{\centering{ }}
        \captionsetup{font={small,bf,stretch=1}}
        \label{table1}
    \end{table}


     2 
    \begin{table}[H]
        \centering
        \setlength{\tabcolsep}{4mm}
        \begin{threeparttable}  
            \begin{tabular}{c|ccccc}  
                \hline
                \diagbox{ }{ (s)}{ } & $t=20$ &$t=40$ &$t=60$&$t=70$&$t=80$\\ \hline
                PAVF  & $6.4571$ &$17.8073$ &$31.6309$&$38.1617$&$43.9885$\\ \hline
                AVF  & $19.0332$ &$48.5705$ &$78.6096$&$91.5572$&$122.6046$\\ \hline
            \end{tabular} 
        \end{threeparttable}
        \caption{\centering{ }}
        \captionsetup{font={small,bf,stretch=1}}
        \label{table2}%label caption .
    \end{table}

     3 
    \begin{table}[H]
        \centering
        \setlength{\tabcolsep}{4mm}
        \begin{threeparttable}  
            \begin{tabular}{c|ccccc}  
                \hline
                \diagbox{ \tnote{*}}{ (s)}{ } & $t=20$ &$t=40$ &$t=60$&$t=70$&$t=80$\\ \hline
                PAVF  & $6.4571$ &$17.8073$ &$31.6309$&$38.1617$&$43.9885$\\ \hline
                AVF  & $19.0332$ &$48.5705$ &$78.6096$&$91.5572$&$122.6046$\\ \hline
            \end{tabular} 
            \captionsetup{font={small,bf,stretch=1}}
            \begin{tablenotes}    %   
                \footnotesize               % 
                \item[*]        %       
            \end{tablenotes}            % 
        \end{threeparttable}
        \caption{\centering{ }}
        \label{table3}
    \end{table}
     

     ref $\ref{table1}$ $\ref{table2}$ $\ref{table3}$ 









    \newpage
    \section{ }
     
    % \section{} 
    \markboth{\text{ }\thesection\text{ } \quad  }{\text{ }\thesection\text{ } \quad  }  

    \begin{algorithm}[!ht]
        \renewcommand{\algorithmicrequire}{\textbf{ }}
        \renewcommand{\algorithmicensure}{\textbf{ }}
        \caption{ }
        \label{power0}
        \begin{algorithmic}[1] %  
            \REQUIRE   % input  
            \ENSURE  % output  
            
            \STATE  ;
            \STATE  ; 
            
             % for loop
            \FORALL {$k=1,2,\cdots,Iter$}
                \STATE $N=N(\widehat{\varPhi}^0,\widehat{\varPhi}^1)$;
                \STATE $A=A(\widehat{\varPhi}^1)$;
                \STATE $\widehat{\varPhi}^2=A\backslash N$;
                \STATE $\widehat{\varPhi}^0=\widehat{\varPhi}^1$;
                \STATE $\widehat{\varPhi}^1=\widehat{\varPhi}^2$;
            \ENDFOR
            
             %if ... else
            \IF {$a < 0$}
                \STATE $y = 2$;
            \ELSE
                \STATE $y = 4$;
            \ENDIF
            
            % while 
            \WHILE {a < 0}
                \STATE a--;
            \ENDWHILE
            
             do ... while
            \REPEAT
                \STATE a--
            \UNTIL {a < 0}
            
            \STATE \textbf{return} $\widehat{\varPhi}^1$.
        \end{algorithmic}
    \end{algorithm}

     for,while,if,do......while 
    










    \newpage
    %\vspace*{0pt}
    \section{ }
    % \section{} 
    \markboth{\text{ }\thesection\text{ } \quad  }{\text{ }\thesection\text{ } \quad  }  

     

    \begin{enumerate}[1)]
    \item  word pdf \textcolor{red}{ cover.pdf main.tex } latex 
    \item  Untitled\_dedai.m main.tex mcode matlab python,c,java 
    \item  href pdf bib.tex, 
    \item  section markboth section 
    \end{enumerate}






    \newpage
    \section*{ }
    \phantomsection
    \addcontentsline{toc}{section}{ }
    \fancyhead[C]{\zihao{5}\songti{ }}
     
    \lstset{
        caption={ },
        aboveskip=0pt,% 
    }
    \lstinputlisting[language=Matlab]{Untitled_dedai.m}

    \textcolor{red}{ }



    \newpage
    \section*{ }
    \phantomsection
    \addcontentsline{toc}{section}{ }
    \fancyhead[C]{\zihao{5}\songti{ }}
    \begingroup  %  thebibliography 
    \renewcommand{\section}[2]{} 

    \begin{thebibliography}{99}
        %\addtolength{\itemsep}{-1pt}  %  
        \addtolength{\itemsep}{0em} %  
        \setlength{\itemsep}{0pt}
        \bibitem{ref1}\href{https://www.sciencedirect.com/science/article/pii/S0893965921003098}{H.G. Lee, A non-iterative and unconditionally energy stable method the Swift--Hohenberg equation with quadratic--cubic nonlinearity[J]. Applied Mathematics Letters, 2021, 123: 107579.}
        \bibitem{ref2}\href{https://www.sciencedirect.com/science/article/pii/S0893965921000847}{Z.F. Weng, Y.F. Deng, Q.Q. Zhuang, S.Y. Zhai, A fast and efficient numerical algorithm for Swift--Hohenberg equation with a nonlocal nonlinearity[J]. Applied Mathematics Letters, 2021, 118: 107170.}

    \end{thebibliography}
    
    \endgroup




    \newpage
    \section*{ }
    \phantomsection
    \addcontentsline{toc}{section}{ }
    \fancyhead[C]{\zihao{5}\songti{ }}
     
     

    2023 2023 05 12
\end{document}